\chapter{Conclusion}
\label{ch:conclusion}

This chapter summarises what was built and what was learned, states the limitations
of the work honestly, and sets out the improvements that follow most directly from
them.

\section{Summary}
\label{sec:concl-summary}

LectureAssist is an end-to-end, fully local agentic system that turns a raw lecture
video or academic document into validated, difficulty-controlled study and assessment
material. It fuses two extraction channels --- scene-adaptive OCR that classifies each
frame before preprocessing it, so that slides, blackboards, whiteboards and on-screen
text are each handled correctly, and Faster-Whisper speech transcription --- into a
single timestamped lecture timeline. That timeline is summarised, chunked, embedded
and indexed, and then consumed by a multi-agent pipeline that plans questions across
topics, retrieves its own evidence when it judges its context insufficient, generates
one question at a time, validates each against difficulty, grounding and instruction
compliance, and, on failure, chooses \emph{how} to remediate before retrying. Around
this sit a retrieval-grounded chat interface, a multi-tier grader that falls back from
exact matching through fuzzy similarity to LLM rubric evaluation, and an adaptive
engine that steers subsequent quizzes toward the topics a student is weakest on.
Everything runs on open-weight Gemma~3 models served locally through Ollama on a
single free-tier GPU, with no paid API anywhere in the system.

The project's second contribution is an evaluation that takes seriously a question the
field usually does not ask: not \emph{does the generator write good questions}, but
\emph{does it mark the right answer}. Question generation and answer generation are
treated as two distinct agents, and their answers are cross-checked by having a
larger model re-solve every question without seeing the marked option. Measured this
way against a plain single-shot baseline on a college calculus chapter, the agentic
pipeline wins on judged quality at every difficulty and wins decisively on answer
correctness where it matters most --- $99\%$ versus $81\%$ judge-verified on hard
questions --- at roughly seven times the wall-clock cost.

The most instructive result, however, is a negative one. Judge scores turned out to be
compressed near the ceiling and to discriminate poorly between pipelines, while the
duplicate rate discriminated sharply; and at hard difficulty the agentic pipeline
collapses onto a narrow slice of the chapter and repeats itself so badly that a
repetition-aware metric ranks the plain baseline \emph{above} it. A quality-only
evaluation --- the kind most commonly reported --- would have shown the agentic
pipeline winning cleanly across the board and would have concealed this entirely.

\section{Limitations}
\label{sec:concl-limits}

\paragraph{The judge and the Answer Generator are LLMs, not humans.} All quality
scores are produced by \texttt{gemma3:12b} rather than by human raters. LLM judges are
known to reward fluency and verbosity, and the judge here shares a model family with
the generator, so it plausibly shares its blind spots. Absolute scores should be read
as a relative signal between pipelines, not as an objective measure of question
quality.

\paragraph{The Answer Generator is weak, which bounds every answer-correctness
number.} On the known-answer calibration set the solver scores only $55\%$, and the
generator model $70\%$. Generator--solver agreement of $88$--$97\%$ therefore cannot be
interpreted as that fraction of questions being verifiably correct; two models from
the same family agree partly because they behave alike. This is why the calibration is
reported --- but it means the answer-correctness results establish a ranking between
pipelines, not a trustworthy absolute correctness rate. A stronger or
non-Gemma solver would tighten this considerably.

\paragraph{The CoT and PoT comparison is not yet complete.} The four-pipeline run is
implemented and the prompts are fixed, but the completed measurements reported in
Chapter~\ref{ch:results} cover the agentic versus plain non-agentic comparison; the
CoT and PoT rows are shown as pending rather than filled with estimates.

\paragraph{Single domain, single document.} All results come from one chapter of one
subject (limits, from OpenStax \emph{Calculus, Volume~1}). Mathematics is unusually
favourable to program-of-thought prompting and unusually hostile to OCR, and nothing
here establishes that the findings transfer to a humanities lecture or to a
lab-based science course.

\paragraph{The extraction layer bounds everything above it.} OCR on handwritten
board work is imperfect, transcription degrades with accent and background noise, and
mathematics rendered as an image in a PDF is not extracted at all. Every defect in
extraction propagates silently upward: a missing board block looks exactly like a
board that was never written on, and the question generator has no way to know that
material is absent.

\paragraph{The agentic pipeline's diversity failure at hard difficulty is unresolved.}
It is measured, diagnosed --- the planner and retriever converge on the narrow region
of the chapter that supports difficult questions, and each retry draws from that same
region again --- and reported, but not yet fixed.

\paragraph{Deployment is not production-grade.} The system runs on a free Colab
runtime behind an ngrok tunnel, with a single-user interface, no authentication, and
session state on Google Drive. This is adequate to demonstrate and evaluate the
system; it is not a deployable service.

\section{Future Improvement}
\label{sec:concl-future}

\paragraph{Fix the diversity collapse first.} It is the clearest defect the evaluation
exposed and the one that currently negates the agentic pipeline's advantage on hard
questions. Two mechanisms follow directly from the diagnosis: maintain a coverage map
of which regions of the document have already been used and penalise the planner for
returning to them, and make the generator aware of the questions already produced in
the current run so that a retry is required to be \emph{different}, not merely
\emph{valid}. A duplicate check could also be moved inside the validation loop, so
that a repetitive question is rejected at generation time rather than discovered at
evaluation time.

\paragraph{Strengthen answer verification.} The weakest link in the evaluation is the
solver. Verifying computational answers symbolically with SymPy --- as is already done
when constructing the calibration set --- would replace an LLM's opinion with a
machine-checked fact for exactly the question type where the generator is most likely
to be wrong. Using a solver from a different model family, or an ensemble, would
break the shared-blind-spot problem that currently inflates the agreement rate.

\paragraph{Human evaluation.} A modest human study --- a few instructors rating a
sample of generated questions, and independently answering them --- would anchor the
LLM judge to a ground truth and quantify how far its scores can be trusted. This is
the single highest-value addition to the evaluation.

\paragraph{Broaden the domain.} Repeating the comparison on a non-mathematical subject
would test whether the agentic advantage in answer correctness is general or is an
artifact of a domain in which answers are computed and therefore easy to get wrong.

\paragraph{Complete and extend the pipeline comparison.} Finishing the CoT and PoT
runs will answer whether a better prompt can recover most of the agentic pipeline's
benefit at a fraction of its cost --- a question with direct practical consequences,
given the seven-fold cost difference. A hybrid worth testing is a CoT or PoT generator
placed \emph{inside} the agentic validate-and-retry loop, combining derived answers
with external validation.

\paragraph{Productionisation.} Multi-user support with authentication, a persistent
server rather than an ephemeral notebook runtime, and quantised models for lower-VRAM
consumer GPUs would move the system from a demonstrated prototype toward something a
class of students could actually use.
